\documentclass[11pt]{article}
\usepackage[margin=25mm]{geometry}
\usepackage{fontspec,unicode-math}
\setmainfont{Latin Modern Roman}
\setmathfont{Latin Modern Math}
\usepackage{amsmath,amssymb,mathtools}
\usepackage{xurl,hyperref}
\hypersetup{colorlinks=true,urlcolor=blue,linkcolor=blue}
\setcounter{secnumdepth}{0}
\setlength{\parindent}{0pt}
\setlength{\parskip}{5pt}
\setlength{\emergencystretch}{3em}
\newcommand{\sref}[2]{\hyperlink{source:#1:#2}{[S#2]}}
\newcommand{\cataloguescope}[1]{\paragraph{Source question.}#1}
\newcommand{\floor}[1]{\left\lfloor#1\right\rfloor}
\begin{document}
\section{Adjoint-Free Approximation of Nonsymmetric Low-Rank Matrices}\label{statement}

\subsection{Definitions and mathematical statement}

\paragraph{Shared notation.}Standard number systems are denoted by $\mathbb N,\mathbb Z,\mathbb Q,\mathbb R,\mathbb C$; index ranges, computational models and specialized notation follow the cited problem source and the definitions below.


The adjoint $A^*$ satisfies $\langle Ax,y\rangle=\langle x,A^*y\rangle$. A forward sketch queries $A$ on selected vectors; adjoint-free access forbids corresponding queries to $A^*$. Low-rank approximation minimizes the stated matrix or operator error.

The rank of a matrix is the dimension of its column space. Matrix completion reconstructs a low-rank matrix from a subset of its entries; matrix sensing reconstructs it from specified linear measurements. Tensor decomposition represents a multiway array as a sum of simple outer products.
\paragraph{Statement recorded in the supplied source.}
Determine when and with what sample complexity a nonsymmetric low-rank matrix or nonselfadjoint compact operator can be approximated using only forward sketches, without querying its adjoint. \sref{statement}{2}
\paragraph{Scope qualification.}
The operator class, sampling distribution, error norm, and target success probability require the source.
\paragraph{Source status.}
The supplied catalogue labels this record partially solved. Its source-status metadata was checked against the supplied snapshot on 2026-10-11. This import does not independently establish current open status or mathematical correctness. The catalogue literature assessment is dated 2026-08-22.
\subsection{Short English statement}
Approximate nonsymmetric low-rank objects using forward access alone and determine the required sample count.
\subsection{Sources}
\begin{itemize}
\item[\hypertarget{source:statement:1}{S1}] Ulam AI and the UnsolvedMath Contributors, UnsolvedMath: A Curated Collection of Open Mathematics Problems, catalogue entry. Catalogue curation and metadata: CC BY 4.0. The supplied source question is retained; definitions and the English summary are editorial additions.
\url{https://huggingface.co/datasets/ulamai/UnsolvedMath}.
\item[\hypertarget{source:statement:2}{S2}] Deep Learning for PDE-based Inverse Problems (2025), Oberwolfach Reports, DOI 10.4171/owr/2024/48, pp. 2805--2900
\url{https://doi.org/10.4171/owr/2024/48}.
\item[\hypertarget{source:statement:3}{S3}] Boullé et al., Operator learning without the adjoint
\url{https://www.jmlr.org/papers/v25/24-0162.html}.
\textit{Underlying Oberwolfach Reports material retains its source-specific licence, including CC BY-SA 4.0 where applicable; see THIRD\_PARTY\_NOTICES.md.}
\end{itemize}
\paragraph{Preparation and review.}
Original question attributed to the cited source and UnsolvedMath contributors. Codex (GPT-6; exact serving revision not exposed) assisted with definitions, English summaries, source inspection and formatting. No new human mathematical authorship or independent proof verification is asserted. Catalogue identity checked against the complete supplied registry and existing repository statements; known flagged repetitions and obvious identical conjectures excluded. Full mathematical-equivalence/literature review has not been performed. Statement preparation and final review followed contribution prompts C and D; automated checks verify structure and identity, not open status.
\label{end:statement}
\end{document}
